This appendix collects the full round-level tables behind
\figref{fig:r5_debate_dynamics_label_based} and
\figref{fig:r5_debate_dynamics_label_free}. The tables expand the mechanism
diagnostics used in \secref{sec:experiments} and \secref{sec:supply}. They ask two
questions: whether the selected agents supply correct proposals that vote
misses, and whether the matched readout converts that supply into final
accuracy without excessive damage.

\paragraph{Metrics and Reading Guide.}
Vote is the majority-vote accuracy on the proposal multiset at a given round.
Proposal Hit is the probability that at least one visible proposal at that
round contains the correct answer. Headroom $H$ is the gap between these two
events: it counts cases where the proposal multiset contains the correct answer
but majority vote misses it. Final is the end-of-run
\texttt{vote\_then\_judge\_tiebreak} accuracy. Gain $G$ is Final minus Vote.
Recovery $R$ is the final accuracy on headroom instances. New Correctness $I$
is the final accuracy on instances where the visible proposal multiset did not
contain the correct answer and Vote was wrong. Damage $L$ is the failure rate
on instances where Vote is already correct. Thus Proposal Hit and Headroom
diagnose proposal supply, while Recovery, New Correctness, and Damage diagnose
readout conversion.

These definitions make the appendix a supply-readout analysis rather than a
second main-results table. Higher Proposal Hit or Headroom means that the
selected agents surface more correct proposals. Higher Recovery means that the
readout converts more surfaced correct minorities into final answers. Lower
Damage means that the readout preserves cases already solved by majority vote.
New Correctness tracks gains not explained by the proposal-hit supply at the
chosen round. Final accuracy combines these effects, so it can hide whether a
gain comes from better proposal supply, better readout conversion, or lower
damage.

The main pattern is that useful proposal supply appears early and then contracts
as debate moves toward consensus. At R0, all four debate societies have
7--8 points of recoverable headroom on average. By R5, the remaining headroom
falls to about one point. This supports the fixed-initial-proposal view in
\secref{sec:headroom_lvd}: the society often surfaces correct minorities before
interaction, and later rounds can reduce the set of alternatives even when vote
accuracy improves. Debate gains therefore depend on recovering useful minority
answers before proposal-hit coverage collapses.

The recovery and damage columns show why higher proposal-hit supply is not
enough by itself. For example, AC-Greedy Debate has the highest average R0
proposal-hit rate, 84.93\%, but its final lift over the R0 vote is 1.64 points.
SC-Greedy Debate has lower average R0 proposal-hit rate, 83.88\%, but larger
average lift, 2.48 points, because its starting vote is lower and more of its
gain comes from converting missed proposals. These differences are why the main
text reports headroom, recovery, and damage rather than final accuracy alone.

\subsubsection{Average Recovery Diagnostics}

\begin{table}[t]
\centering
\caption{Round-level recovery diagnostics averaged over the five benchmarks.
For each round, Vote and Hit are the round-specific majority-vote and
proposal-hit accuracies, while Final is the end-of-run
\texttt{vote\_then\_judge\_tiebreak} accuracy. We define $H=P-V$, $G=D-V$,
$\rho=G/H$, $R=\Pr(D=1\mid P=1,V=0)$,
$I=\Pr(D=1\mid P=0,V=0)$, and $L=\Pr(D=0\mid V=1)$, where $D$ is the final
decision. All values except $\rho$ are percentages; averages are unweighted
over benchmarks.}
\scriptsize
\resizebox{\textwidth}{!}{%
\begin{tabular}{llrrrrrrrrr}
\toprule
Method & Round & Vote & Hit & Final & $H$ & $G$ & $\rho$ & $R$ & $I$ & $L$ \\
\midrule
Vanilla Debate & R0 & 75.65 & 83.03 & 77.92 & 7.39 & 2.27 & 0.31 & 53.49 & 9.94 & 5.00 \\
 & R1 & 77.85 & 85.05 & 77.92 & 7.20 & 0.07 & 0.01 & 43.10 & 1.23 & 4.25 \\
 & R2 & 77.84 & 81.86 & 77.92 & 4.02 & 0.08 & 0.02 & 24.33 & 0.14 & 1.58 \\
 & R3 & 78.05 & 80.07 & 77.92 & 2.03 & -0.13 & -0.06 & 43.39 & 0.60 & 1.37 \\
 & R4 & 78.13 & 79.04 & 77.92 & 0.92 & -0.21 & -0.23 & 38.89 & 0.13 & 1.08 \\
 & R5 & 77.99 & 78.98 & 77.92 & 0.99 & -0.08 & -0.08 & 10.00 & 0.13 & 0.33 \\
\midrule
Top-4 Debate & R0 & 76.25 & 84.11 & 78.28 & 7.87 & 2.03 & 0.26 & 40.88 & 7.26 & 4.53 \\
 & R1 & 78.51 & 85.15 & 78.28 & 6.64 & -0.23 & -0.03 & 29.46 & 0.38 & 3.67 \\
 & R2 & 78.72 & 82.22 & 78.28 & 3.51 & -0.44 & -0.13 & 28.94 & 0.48 & 2.16 \\
 & R3 & 78.30 & 80.93 & 78.28 & 2.63 & -0.02 & -0.01 & 30.09 & 0.00 & 1.18 \\
 & R4 & 78.11 & 79.66 & 78.28 & 1.56 & 0.17 & 0.11 & 42.93 & 0.00 & 0.73 \\
 & R5 & 78.16 & 79.27 & 78.28 & 1.11 & 0.11 & 0.10 & 17.78 & 0.00 & 0.25 \\
\midrule
AC-Greedy Debate & R0 & 76.84 & 84.93 & 78.48 & 8.09 & 1.64 & 0.20 & 48.94 & 7.70 & 6.01 \\
 & R1 & 78.63 & 85.72 & 78.48 & 7.09 & -0.15 & -0.02 & 34.76 & 1.68 & 4.23 \\
 & R2 & 78.43 & 82.49 & 78.48 & 4.06 & 0.06 & 0.01 & 45.25 & 0.40 & 2.46 \\
 & R3 & 78.73 & 80.84 & 78.48 & 2.10 & -0.25 & -0.12 & 37.24 & 0.00 & 1.58 \\
 & R4 & 78.50 & 79.97 & 78.48 & 1.47 & -0.02 & -0.01 & 39.44 & 0.00 & 0.64 \\
 & R5 & 78.41 & 79.60 & 78.48 & 1.19 & 0.07 & 0.06 & 32.00 & 0.00 & 0.11 \\
\midrule
SC-Greedy Debate & R0 & 75.64 & 83.88 & 78.12 & 8.23 & 2.48 & 0.30 & 48.52 & 7.20 & 5.34 \\
 & R1 & 77.37 & 85.35 & 78.12 & 7.98 & 0.75 & 0.09 & 34.76 & 1.45 & 3.55 \\
 & R2 & 78.32 & 82.25 & 78.12 & 3.93 & -0.20 & -0.05 & 25.29 & 0.00 & 2.58 \\
 & R3 & 78.33 & 80.42 & 78.12 & 2.09 & -0.21 & -0.10 & 29.56 & 0.00 & 1.37 \\
 & R4 & 77.69 & 79.65 & 78.12 & 1.97 & 0.43 & 0.22 & 24.71 & 0.00 & 0.42 \\
 & R5 & 78.22 & 79.19 & 78.12 & 0.97 & -0.10 & -0.10 & 6.67 & 0.00 & 0.26 \\
\bottomrule
\end{tabular}
}
\label{tab:round_recovery_diagnostics}
\end{table}

\paragraph{Interpretation.}
\tabref{tab:round_recovery_diagnostics} separates proposal supply from readout
conversion. Initial headroom is substantial for every method, which means the
selected societies surface correct answers that majority vote does not extract.
However, recovery falls sharply after the first few rounds because debate
compresses the proposal set. Damage also falls over rounds, since fewer
vote-correct cases are changed once the agents converge. New Correctness is
largest at R0 and becomes small after the debate trace has mostly converged,
which means most final gains are explained by surfaced proposal supply rather
than late invention of new correct answers. The tradeoff is
therefore not monotone: later rounds can make the majority more stable, but
they can also remove the minority alternatives that the readout needs to
recover.

\subsubsection{Five-Round New-Seed Debate Results}

\begin{table}[t]
\centering
\caption{Five-round debate results for the new seed teams. Final is accuracy
after \texttt{vote\_then\_judge\_tiebreak}; Lift is Final minus Init Vote.
Bold marks the best final accuracy within each dataset. Average rows report the
unweighted mean across the five benchmarks.}
\small
\resizebox{\textwidth}{!}{%
\begin{tabular}{llrrrrrrrr}
\toprule
Dataset & Method & N & Init Vote & R5 Vote & Final & Init Hit & R5 Hit & Lift & Judge Ties \\
\midrule
MATH500 & Vanilla Debate & 300 & 91.33 & 91.33 & 91.33 & 93.33 & 91.67 & +0.00 & 1 \\
MATH500 & AC-Greedy Debate & 300 & 91.33 & 91.33 & 91.67 & 94.33 & 91.67 & +0.33 & 2 \\
MATH500 & SC-Greedy Debate & 300 & 92.00 & 92.00 & \textbf{92.00} & 93.67 & 92.67 & +0.00 & 0 \\
OlympiadBench & Vanilla Debate & 474 & 80.38 & 81.65 & 81.22 & 83.54 & 82.28 & +0.84 & 7 \\
OlympiadBench & AC-Greedy Debate & 474 & 79.11 & 81.43 & \textbf{81.43} & 82.91 & 81.86 & +2.32 & 7 \\
OlympiadBench & SC-Greedy Debate & 474 & 80.59 & 80.59 & 80.80 & 83.33 & 81.22 & +0.21 & 5 \\
AMC12 & Vanilla Debate & 739 & 91.07 & 92.96 & \textbf{92.96} & 94.32 & 93.37 & +1.89 & 4 \\
AMC12 & AC-Greedy Debate & 739 & 90.53 & 91.88 & 91.61 & 95.13 & 92.83 & +1.08 & 9 \\
AMC12 & SC-Greedy Debate & 739 & 89.85 & 92.02 & 91.88 & 93.91 & 92.56 & +2.03 & 2 \\
GPQA & Vanilla Debate & 346 & 47.11 & 52.02 & 51.73 & 67.63 & 54.91 & +4.62 & 4 \\
GPQA & AC-Greedy Debate & 346 & 48.55 & 51.73 & \textbf{52.02} & 67.63 & 54.62 & +3.47 & 3 \\
GPQA & SC-Greedy Debate & 346 & 47.11 & 51.16 & 50.58 & 66.47 & 53.18 & +3.47 & 4 \\
MMLU-Redux & Vanilla Debate & 300 & 68.33 & 72.00 & 72.33 & 76.33 & 72.67 & +4.00 & 2 \\
MMLU-Redux & AC-Greedy Debate & 300 & 74.67 & 75.67 & \textbf{75.67} & 84.67 & 77.00 & +1.00 & 1 \\
MMLU-Redux & SC-Greedy Debate & 300 & 68.67 & 75.33 & 75.33 & 82.00 & 76.33 & +6.67 & 2 \\
\midrule
Average & Vanilla Debate & -- & 75.65 & 77.99 & 77.92 & 83.03 & 78.98 & +2.27 & 3.6 \\
Average & AC-Greedy Debate & -- & 76.84 & 78.41 & \textbf{78.48} & 84.93 & 79.60 & +1.64 & 4.4 \\
Average & SC-Greedy Debate & -- & 75.64 & 78.22 & 78.12 & 83.88 & 79.19 & +2.48 & 2.6 \\
\bottomrule
\end{tabular}
}
\label{tab:r5_new_seed_debate}
\end{table}

\paragraph{Interpretation.}
\tabref{tab:r5_new_seed_debate} shows that the same mechanism appears at the
dataset level. On GPQA and MMLU-Redux, initial proposal-hit rates are much
higher than initial vote accuracy, and debate produces the largest lifts. On
MATH500 and AMC12, the initial vote is already high, so the available room for
final improvement is smaller even when proposal-hit coverage is high. This is
consistent with \propref{prop:recovery_damage_decomposition}: final gain is
largest when there is recoverable headroom and the readout converts it without
large damage.

The construction methods also behave differently across datasets. AC-Greedy
Debate gives the best average final accuracy, but SC-Greedy Debate gives the
largest average lift from its own initial vote. This supports the role of
label-free coverage as a practical supply constructor, while also showing why
the supervised construction remains a useful reference.
